% ============================================================
% Paper 45: Ringtheorie und kommutative Algebra (Deutsch)
% Autor: Michael Fuhrmann
% Projekt: Specialist Mathematics, Build 135
% Datum: 2026-03-12
% ============================================================

\documentclass[12pt,a4paper]{amsart}
\usepackage{amsmath,amssymb,amsthm}
\usepackage{mathtools}
\usepackage[utf8]{inputenc}
\usepackage[T1]{fontenc}
\usepackage{hyperref}
\usepackage{enumitem,booktabs,array}
\usepackage{geometry}
\geometry{margin=2.5cm}

% ---- Theorem-Umgebungen ----
\newtheorem{theorem}{Satz}[section]
\newtheorem{lemma}[theorem]{Lemma}
\newtheorem{corollary}[theorem]{Korollar}
\newtheorem{proposition}[theorem]{Proposition}
\newtheorem{conjecture}[theorem]{Vermutung}
\theoremstyle{definition}
\newtheorem{definition}[theorem]{Definition}
\newtheorem{example}[theorem]{Beispiel}
\newtheorem{remark}[theorem]{Bemerkung}
\theoremstyle{definition}
\newtheorem{algorithm}[theorem]{Algorithmus}

% ---- Makros ----
\newcommand{\ggT}{\operatorname{ggT}}

% ---- Operatoren ----
\DeclareMathOperator{\Spec}{Spec}
\DeclareMathOperator{\Ann}{Ann}
\DeclareMathOperator{\rad}{rad}
\DeclareMathOperator{\Nil}{Nil}
\DeclareMathOperator{\Jac}{Jac}
\DeclareMathOperator{\lcm}{kgV}
\DeclareMathOperator{\htOp}{ht}
\DeclareMathOperator{\codim}{codim}
\DeclareMathOperator{\chr}{char}
\DeclareMathOperator{\Frac}{Frac}
\DeclareMathOperator{\Cl}{Cl}

\title{Ringtheorie und kommutative Algebra\\
\large Ideale, Faktorisierung, noethersche Ringe und Gröbner-Basen}

\author{Michael Fuhrmann}
\address{Specialist Mathematics Project, Build 135}
\email{specialist-maths@localhost}
\date{\today}

\begin{document}

\begin{abstract}
Dieses Paper gibt eine in sich geschlossene Einführung in die Ringtheorie und
kommutative Algebra und behandelt die zentralen Objekte und Sätze der modernen
Algebra. Ausgehend von der axiomatischen Definition eines Rings entwickeln wir
die Theorie der Ideale und Quotientenringe, beweisen die Isomorphiesätze und
untersuchen Polynomringe mit ihren Irreduzibilitätskriterien. Die Hierarchie
euklidischer Ring $\Rightarrow$ Hauptidealring (HIR) $\Rightarrow$ faktorieller Ring (UFD)
wird vollständig entwickelt. Anschließend behandeln wir noethersche und artinsche Ringe,
einschließlich des Hilbert-Basissatzes und des Nakayama-Lemmas. Lokalisierung und ihre
Anwendung auf Dedekind-Ringe bilden die Brücke zur algebraischen Zahlentheorie. Das
Paper schließt mit einer Einführung in Gröbner-Basen und den Buchberger-Algorithmus
als Berechnungswerkzeug. Alle Ergebnisse werden mit Beweisen oder Beweissk{}izzen
begleitet und durch konkrete Beispiele illustriert. Die Darstellung stützt sich auf
die Implementierungen in \texttt{ring\_theory.py} und \texttt{commutative\_algebra.py}
des Specialist-Mathematics-Projekts.
\end{abstract}

\maketitle
\tableofcontents

% ============================================================
\section{Einleitung}
% ============================================================

Ringe sind die grundlegenden algebraischen Strukturen, die die Arithmetik der ganzen
Zahlen verallgemeinern. Überall dort, wo Addition und Multiplikation zusammenspielen ---
in der Zahlentheorie, der Geometrie, der Kombinatorik oder der mathematischen Physik ---
liefern ringtheoretische Methoden den vereinheitlichenden Rahmen.

Die historische Entwicklung der Ringtheorie wurde von konkreten Problemen angetrieben.
Gauß untersuchte $\mathbb{Z}[i]$, um seinen Zwei-Quadrate-Satz zu beweisen.
Kummer führte „ideale Zahlen" ein, um die eindeutige Faktorisierung in Kreiskörpern
zu retten. Dedekind formalisierte den modernen Idealbegriff und bewies, dass Ringe
algebraischer ganzer Zahlen genau das sind, was wir heute Dedekind-Ringe nennen ---
Strukturen, in denen eindeutige Faktorisierung auf Idealebene gilt, auch wenn sie auf
Elementebene versagt.

Im zwanzigsten Jahrhundert revolutionierte Emmy Noether~\cite{Noether1921} das Gebiet,
indem sie den Fokus von expliziten Berechnungen auf strukturelle Prinzipien verlagerte.
Ihre aufsteigende Kettenbedingung, die noethersche Eigenschaft, erwies sich als genau
die Hypothese, die benötigt wird, damit Algebra sauber funktioniert. Gleichzeitig zeigte
Hilberts Basissatz, dass Polynomringe diese gute Eigenschaft erben, was das Fundament
der modernen algebraischen Geometrie legte.

Heute bildet die kommutative Algebra --- die Untersuchung kommutativer Ringe und ihrer
Moduln --- den technischen Kern der algebraischen Geometrie, der algebraischen Zahlentheorie
und Teilen der Kombinatorik. Die Sprache von $\Spec$, Lokalisierung und Garben verbindet
reine Algebra mit Geometrie auf eine Weise, die seit Grothendiecks Neuformulierung in
den 1950er und 1960er Jahren entscheidend war.

Dieses Paper deckt folgende Themen ab:
\begin{enumerate}[label=(\roman*)]
  \item Grunddefinitionen und Beispiele von Ringen und Körpern;
  \item Ideale, Quotientenringe und der Chinesische Restsatz;
  \item Ringhomomorphismen und die Isomorphiesätze;
  \item Polynomringe, Gaußsches Lemma und Eisenstein-Kriterium;
  \item Die Hierarchie UFD $\supset$ HIR $\supset$ euklidischer Ring;
  \item Noethersche und artinsche Ringe, Hilbert-Basissatz, Nakayama-Lemma;
  \item Lokalisierung und Dedekind-Ringe;
  \item Gröbner-Basen und der Buchberger-Algorithmus.
\end{enumerate}

\medskip

\noindent\textbf{Notation.} Alle Ringe haben eine multiplikative Eins $1$, und
Ringhomomorphismen bilden $1$ auf $1$ ab. Wir schreiben $\mathbb{Z}$, $\mathbb{Q}$,
$\mathbb{R}$, $\mathbb{C}$ für die ganzen, rationalen, reellen und komplexen Zahlen,
sowie $\mathbb{F}_p = \mathbb{Z}/p\mathbb{Z}$ für den Körper mit $p$ Elementen.

% ============================================================
\section{Grundlagen der Ringtheorie}
% ============================================================

\begin{definition}[Ring]
Ein \emph{Ring} ist eine Menge $R$ zusammen mit zwei binären Operationen $+$
(Addition) und $\cdot$ (Multiplikation), die folgendes erfüllen:
\begin{enumerate}[label=(\roman*)]
  \item $(R,+)$ ist eine abelsche Gruppe mit neutralem Element $0$;
  \item $(R,\cdot)$ ist ein Monoid mit neutralem Element $1$;
  \item Die Multiplikation ist distributiv über die Addition:
        $a(b+c) = ab + ac$ und $(a+b)c = ac + bc$ für alle $a,b,c \in R$.
\end{enumerate}
Der Ring heißt \emph{kommutativ}, wenn $ab = ba$ für alle $a,b \in R$ gilt.
\end{definition}

\begin{example}\label{bsp:ringe}
Die wichtigsten Beispiele sind:
\begin{enumerate}[label=(\alph*)]
  \item $\mathbb{Z}$, $\mathbb{Q}$, $\mathbb{R}$, $\mathbb{C}$ mit der gewöhnlichen
        Addition und Multiplikation;
  \item $\mathbb{Z}/n\mathbb{Z}$ (ganze Zahlen modulo $n$) für beliebiges $n \geq 1$;
  \item Der Polynomring $R[x_1,\ldots,x_n]$ über einem beliebigen Ring $R$;
  \item Der Matrizenring $M_n(R)$ --- nicht kommutativ für $n \geq 2$;
  \item Der Ring der stetigen Funktionen $C([0,1])$ mit punktweisen Operationen;
  \item Quadratische Ganzzahlringe $\mathbb{Z}[\sqrt{d}] = \{a + b\sqrt{d} \mid a,b \in \mathbb{Z}\}$
        für quadratfreies $d \in \mathbb{Z}$.
\end{enumerate}
\end{example}

\begin{definition}[Spezielle Ringklassen]
Sei $R$ ein kommutativer Ring mit $1 \neq 0$.
\begin{itemize}
  \item $R$ ist ein \emph{Integritätsbereich} (oder \emph{nullteilerfreier Ring}),
        wenn aus $ab = 0$ folgt: $a = 0$ oder $b = 0$.
  \item $R$ ist ein \emph{Körper}, wenn jedes von Null verschiedene Element ein
        multiplikatives Inverses besitzt.
  \item Die \emph{Einheitengruppe} ist $R^\times = \{u \in R \mid \exists v: uv = 1\}$.
  \item Die \emph{Charakteristik} $\chr(R)$ ist das kleinste positive $n$ mit
        $n \cdot 1 = 0$, oder $0$, wenn kein solches $n$ existiert.
\end{itemize}
\end{definition}

\begin{proposition}\label{prop:koerper-int}
Jeder Körper ist ein Integritätsbereich. Umgekehrt ist jeder endliche Integritätsbereich
ein Körper.
\end{proposition}

\begin{proof}
Sei $F$ ein Körper und $ab = 0$ mit $a \neq 0$. Dann gilt $b = a^{-1}(ab) = 0$.
Für die Umkehrung sei $R$ ein endlicher Integritätsbereich und $a \in R \setminus \{0\}$.
Die Abbildung $x \mapsto ax$ ist injektiv (da $R$ nullteilerfrei ist) und wegen der
Endlichkeit von $R$ auch surjektiv. Insbesondere hat $ax = 1$ eine Lösung.
\end{proof}

\begin{example}[Einheiten und Nullteiler]
In $\mathbb{Z}/12\mathbb{Z}$: Die Einheiten sind $\{1,5,7,11\}$ (zu $12$ teilerfremde
Elemente), und $2 \cdot 6 = 0$, sodass $2$ ein Nullteiler ist.
In $\mathbb{Z}[i]$ (Gaußsche ganze Zahlen): $(\mathbb{Z}[i])^\times = \{\pm 1, \pm i\}$.
\end{example}

\begin{theorem}[Wedderburns kleiner Satz]
Jeder endliche Schiefkörper ist kommutativ (also ein Körper).
\end{theorem}

\begin{proof}[Beweissk{}izze]
Man verwendet die Klassengleichung der multiplikativen Gruppe und Eigenschaften
von Kreisteilungspolynomen, um zu zeigen, dass das Zentrum des Schiefkörpers
mit dem ganzen Ring übereinstimmen muss. Siehe \cite[Satz~13.1]{Lang2002}.
\end{proof}

% ============================================================
\section{Ideale und Quotientenringe}
% ============================================================

\begin{definition}[Ideal]
Eine Teilmenge $I \subseteq R$ ist ein \emph{(zweiseitiges) Ideal}, wenn:
\begin{enumerate}[label=(\roman*)]
  \item $(I,+)$ eine Untergruppe von $(R,+)$ ist;
  \item $rI \subseteq I$ und $Ir \subseteq I$ für alle $r \in R$ gilt
        (im kommutativen Fall fallen diese beiden Bedingungen zusammen).
\end{enumerate}
Das von $a_1,\ldots,a_n \in R$ \emph{erzeugte Ideal} ist
\[
  (a_1,\ldots,a_n) = \Bigl\{\sum_{i=1}^{n} r_i a_i \,\Big|\, r_i \in R\Bigr\}.
\]
Ein Ideal der Form $(a) = aR$ heißt \emph{Hauptideal}.
\end{definition}

\begin{definition}[Primideale und maximale Ideale]
Sei $R$ ein kommutativer Ring und $\mathfrak{p}, \mathfrak{m} \subsetneq R$ Ideale.
\begin{itemize}
  \item $\mathfrak{p}$ ist ein \emph{Primideal}, wenn aus $ab \in \mathfrak{p}$ folgt:
        $a \in \mathfrak{p}$ oder $b \in \mathfrak{p}$.
  \item $\mathfrak{m}$ ist ein \emph{maximales Ideal}, wenn kein Ideal $I$ mit
        $\mathfrak{m} \subsetneq I \subsetneq R$ existiert.
\end{itemize}
\end{definition}

\begin{theorem}[Charakterisierung über Quotientenringe]
Sei $I$ ein Ideal in einem kommutativen Ring $R$.
\begin{enumerate}[label=(\roman*)]
  \item $I$ ist prim genau dann, wenn $R/I$ ein Integritätsbereich ist.
  \item $I$ ist maximal genau dann, wenn $R/I$ ein Körper ist.
  \item Jedes maximale Ideal ist prim (da jeder Körper ein Integritätsbereich ist).
\end{enumerate}
\end{theorem}

\begin{proof}
(i) $ab + I = (a+I)(b+I) = 0$ in $R/I$ genau dann, wenn $ab \in I$. Die Nullteilerfreiheit
in $R/I$ bedeutet genau, dass daraus $a \in I$ oder $b \in I$ folgt.
(ii) $R/I$ ist ein Körper genau dann, wenn jedes von Null verschiedene Element
$a + I$ eine Einheit ist, d.h. es gibt $b$ mit $ab - 1 \in I$. Dies ist äquivalent
zur Maximalität von $I$: für jedes $r \notin I$ liefert $(I,r) = R$ ein $ar + j = 1$
mit $a \in R$, $j \in I$.
\end{proof}

\begin{example}\label{bsp:ideale-Z}
In $\mathbb{Z}$: Die Primideale sind $(0)$ und $(p)$ für Primzahlen $p$.
Die maximalen Ideale sind genau die $(p)$. Das Ideal $(6) = (2)(3)$ ist weder prim
noch maximal: $\mathbb{Z}/(6) \cong \mathbb{Z}/2 \times \mathbb{Z}/3$ ist kein
Integritätsbereich.
\end{example}

\begin{theorem}[Chinesischer Restsatz]\label{satz:crt}
Sei $R$ ein kommutativer Ring und $I_1,\ldots,I_k$ paarweise koprime Ideale
(d.h. $I_i + I_j = R$ für $i \neq j$). Dann ist die natürliche Abbildung
\[
  R \,\Big/\, \bigcap_{i=1}^k I_i \;\xrightarrow{\;\sim\;}\; \prod_{i=1}^k R/I_i
\]
ein Ringisomorphismus.
\end{theorem}

\begin{proof}
Der Homomorphismus $\varphi\colon R \to \prod R/I_i$, $r \mapsto (r + I_i)_i$, hat
den Kern $\bigcap I_i$. Die Surjektivität (und damit die Bijektivität nach dem Übergang
zum Quotienten) folgt aus der Koprimheitsbedingung durch ein Induktionsargument:
Für jedes $j$ gibt es Elemente $e_j \in R$ mit $e_j \equiv 1 \pmod{I_j}$ und
$e_j \equiv 0 \pmod{I_i}$ für $i \neq j$. Diese ``Idempotenten'' erlauben es,
jedes Ziel $(a_1,\ldots,a_k)$ durch $r = \sum_j a_j e_j$ zu erreichen.
\end{proof}

\begin{corollary}
$\mathbb{Z}/n\mathbb{Z} \cong \prod_{p^k \| n} \mathbb{Z}/p^k\mathbb{Z}$
(Produkt über die Primzahlpotenzen in der Primfaktorzerlegung von $n$).
\end{corollary}

Das \emph{Primspektrum} $\Spec(R)$ ist die Menge aller Primideale von $R$,
versehen mit der \emph{Zariski-Topologie}: Abgeschlossene Mengen haben die Form
$V(I) = \{\mathfrak{p} \in \Spec(R) \mid I \subseteq \mathfrak{p}\}$.
Für $R = \mathbb{Z}/n\mathbb{Z}$ entsprechen die Primideale genau den Primteilern
von $n$ (plus dem Nullideal wenn $n = p$ prim ist).

% ============================================================
\section{Ringhomomorphismen und Isomorphiesätze}
% ============================================================

\begin{definition}[Ringhomomorphismus]
Eine Abbildung $\varphi\colon R \to S$ zwischen Ringen ist ein \emph{Ringhomomorphismus},
wenn
\[
  \varphi(a+b) = \varphi(a) + \varphi(b),\quad
  \varphi(ab) = \varphi(a)\varphi(b),\quad
  \varphi(1_R) = 1_S.
\]
Der \emph{Kern} ist $\ker\varphi = \{r \in R \mid \varphi(r) = 0\}$, ein Ideal von $R$.
Das \emph{Bild} ist $\operatorname{Bild}(\varphi)$, ein Unterring von $S$.
\end{definition}

\begin{theorem}[Erster Isomorphiesatz für Ringe]
Sei $\varphi\colon R \to S$ ein Ringhomomorphismus. Dann gilt
\[
  R/\ker\varphi \;\cong\; \operatorname{Bild}(\varphi).
\]
\end{theorem}

\begin{proof}
Die Abbildung $\bar\varphi\colon R/\ker\varphi \to \operatorname{Bild}(\varphi)$,
definiert durch $\bar\varphi(r + \ker\varphi) = \varphi(r)$, ist wohldefiniert
(unabhängig vom Repräsentanten), ein Ringhomomorphismus, injektiv
(da $\ker\bar\varphi = 0$) und surjektiv auf $\operatorname{Bild}(\varphi)$.
\end{proof}

\begin{theorem}[Zweiter und dritter Isomorphiesatz]
Sei $R$ ein Ring, $I \subseteq J$ Ideale von $R$.
\begin{enumerate}[label=(\roman*)]
  \item \textbf{Zweiter Isomorphiesatz:} Ist $S$ ein Unterring und $I$ ein Ideal,
        so ist $S + I$ ein Unterring, $S \cap I$ ein Ideal von $S$, und
        $S/(S \cap I) \cong (S+I)/I$.
  \item \textbf{Dritter Isomorphiesatz:} $J/I$ ist ein Ideal von $R/I$ und
        $(R/I)/(J/I) \cong R/J$.
  \item \textbf{Korrespondenzsatz:} Es gibt eine Bijektion zwischen den Idealen
        von $R/I$ und den Idealen von $R$, die $I$ enthalten.
\end{enumerate}
\end{theorem}

\begin{example}
Der Auswertungshomomorphismus $\varphi\colon \mathbb{R}[x] \to \mathbb{C}$,
$\varphi(f) = f(i)$, hat den Kern $(x^2 + 1)$ und das Bild $\mathbb{C}$, was
$\mathbb{R}[x]/(x^2+1) \cong \mathbb{C}$ liefert.
\end{example}

\begin{definition}[Nilradikal und Jacobson-Radikal]
Für einen kommutativen Ring $R$:
\begin{itemize}
  \item Das \emph{Nilradikal} ist $\Nil(R) = \{r \in R \mid r^n = 0 \text{ für ein } n\} = \bigcap_{\mathfrak{p} \in \Spec(R)} \mathfrak{p}$.
  \item Das \emph{Jacobson-Radikal} ist $\Jac(R) = \bigcap_{\mathfrak{m} \text{ maximal}} \mathfrak{m}$.
\end{itemize}
\end{definition}

% ============================================================
\section{Polynomringe}
% ============================================================

\begin{definition}[Polynomring]
Für einen kommutativen Ring $R$ besteht der \emph{Polynomring} $R[x]$ aus allen
formalen Summen $f = \sum_{k=0}^n a_k x^k$ mit $a_k \in R$, mit der üblichen
Addition und Faltungsmultiplikation. Allgemeiner bezeichnet $R[x_1,\ldots,x_n]$
den Polynomring in $n$ Variablen.
\end{definition}

\begin{theorem}[Gaußsches Lemma]
Sei $R$ ein faktorieller Ring (UFD) mit Quotientenkörper $K = \Frac(R)$.
Ein Polynom $f \in R[x]$ vom Grad $\geq 1$ ist genau dann irreduzibel in $R[x]$,
wenn $f$ \emph{primitiv} ist (der ggT seiner Koeffizienten ist $1$) und in $K[x]$
irreduzibel ist. Folgerung: $R[x]$ ist ein UFD, wenn $R$ es ist.
\end{theorem}

\begin{proof}
Der Schlüsselschritt ist die ``Gaußsche Produktformel'': Wenn $f, g \in R[x]$
primitiv sind, ist auch $fg$ primitiv. Dies wird bewiesen, indem man annimmt,
dies sei falsch, und modulo einem Primteiler eines Koeffizienten von $fg$ reduziert,
um in $(R/(p))[x]$ (ein Integritätsbereich) einen Widerspruch zu erhalten. Daraus
folgt die Multiplikativität des Inhalts, und die Äquivalenz der Irreduzibilität in
$R[x]$ und $K[x]$ ergibt sich. Siehe \cite[IV, \S2]{Lang2002}.
\end{proof}

\begin{theorem}[Eisenstein-Kriterium]\label{satz:eisenstein}
Sei $R$ ein UFD, $\mathfrak{p}$ ein Primideal, und
$f = a_n x^n + \cdots + a_1 x + a_0 \in R[x]$ mit $n \geq 1$.
Angenommen:
\begin{enumerate}[label=(\roman*)]
  \item $a_n \notin \mathfrak{p}$,
  \item $a_0, a_1, \ldots, a_{n-1} \in \mathfrak{p}$,
  \item $a_0 \notin \mathfrak{p}^2$.
\end{enumerate}
Dann ist $f$ irreduzibel in $\Frac(R)[x]$.
\end{theorem}

\begin{proof}
Angenommen, $f = gh$ mit $g, h \in R[x]$ vom positiven Grad. Reduktion modulo $\mathfrak{p}$
liefert $\bar{f} = a_n x^n = \bar{g}\bar{h}$ in $(R/\mathfrak{p})[x]$, einem Integritätsbereich.
Also gilt $\bar{g} = cx^r$ und $\bar{h} = dx^s$ für $r + s = n$, $r, s \geq 1$.
Dann ist $\mathfrak{p} \mid g(0)$ und $\mathfrak{p} \mid h(0)$, also
$\mathfrak{p}^2 \mid g(0)h(0) = a_0$ --- Widerspruch zu (iii).
\end{proof}

\begin{example}
\begin{enumerate}[label=(\alph*)]
  \item $x^p - p$ ist über $\mathbb{Q}$ irreduzibel nach Eisenstein mit $\mathfrak{p}=(p)$.
  \item Das $p$-te Kreisteilungspolynom $\Phi_p(x) = x^{p-1} + \cdots + x + 1$ ist
        irreduzibel: Substitution $x \mapsto x+1$ und Anwendung von Eisenstein mit $p$.
\end{enumerate}
\end{example}

\begin{theorem}
Für einen kommutativen Ring $R$ gilt für $R[x]$:
\begin{itemize}
  \item Ist $R$ ein Integritätsbereich, so auch $R[x]$.
  \item $\chr(R[x]) = \chr(R)$.
  \item $(R[x])^\times = R^\times$ (für Integritätsbereiche $R$).
  \item $\mathbb{F}_q[x]$ (für $q = p^n$ eine Primpotenz) ist ein euklidischer Ring,
        also HIR und UFD.
\end{itemize}
\end{theorem}

% ============================================================
\section{Eindeutige Faktorisierung: UFD, HIR und euklidische Ringe}
% ============================================================

\begin{definition}[Teilbarkeit und Irreduzible]
In einem Integritätsbereich $R$:
\begin{itemize}
  \item $a$ \emph{teilt} $b$ (geschrieben $a \mid b$), wenn $b = ac$ für ein $c \in R$.
  \item $p \in R$ ist \emph{irreduzibel}, wenn aus $p = ab$ folgt, dass $a$ oder $b$
        eine Einheit ist.
  \item $p \in R$ ist \emph{prim}, wenn $(p)$ ein Primideal ist, d.h. aus $p \mid ab$
        folgt $p \mid a$ oder $p \mid b$.
  \item Jedes prime Element ist irreduzibel; in einem UFD gilt auch die Umkehrung.
\end{itemize}
\end{definition}

\begin{definition}[UFD, HIR, euklidischer Ring]
\begin{itemize}
  \item Ein Integritätsbereich $R$ ist ein \emph{faktorieller Ring} (UFD, unique
        factorization domain), wenn jedes von Null verschiedene Nichteinheit eine
        eindeutige Zerlegung in Irreduzible besitzt (bis auf Reihenfolge und Einheiten).
  \item $R$ ist ein \emph{Hauptidealring} (HIR, principal ideal domain, PID), wenn
        jedes Ideal ein Hauptideal ist.
  \item $R$ ist ein \emph{euklidischer Ring}, wenn eine Funktion
        $\delta\colon R \setminus \{0\} \to \mathbb{N}_0$ existiert mit: Für alle
        $a \in R$, $b \in R \setminus \{0\}$ gibt es $q, r \in R$ mit $a = qb + r$
        und entweder $r = 0$ oder $\delta(r) < \delta(b)$.
\end{itemize}
\end{definition}

\begin{theorem}[Hierarchie der Ringklassen]\label{satz:hierarchie}
Es gelten die Implikationen (keine ist umkehrbar):
\[
  \text{Euklidischer Ring} \;\Rightarrow\; \text{HIR} \;\Rightarrow\; \text{UFD}
  \;\Rightarrow\; \text{Integritätsbereich}.
\]
\end{theorem}

\begin{proof}[Beweissk{}izze]
\textit{Euklidisch $\Rightarrow$ HIR:} Sei $I \neq 0$ ein Ideal und $b \in I$ mit
minimalem $\delta(b)$. Für beliebiges $a \in I$ schreibe $a = qb + r$. Da $r = a - qb
\in I$ und $\delta(b)$ minimal ist, muss $r = 0$ gelten, also $a \in (b)$.

\textit{HIR $\Rightarrow$ UFD:} In einem HIR ist jedes irreduzible Element prim
(da $(p)$ maximal, also prim ist). Existenz von Zerlegungen folgt aus der aufsteigenden
Kettenbedingung (alle HIR sind noethersch). Eindeutigkeit folgt aus der Primeigenschaft.

\textit{UFD $\Rightarrow$ Integritätsbereich:} Per Definition.

Gegenbeispiele für die Umkehrungen: $\mathbb{Z}[\sqrt{-5}]$ ist kein UFD:
$6 = 2 \cdot 3 = (1+\sqrt{-5})(1-\sqrt{-5})$ --- zwei verschiedene Zerlegungen.
$\mathbb{Z}[x]$ ist ein UFD, aber kein HIR: Das Ideal $(2,x)$ ist nicht prim.
$k[x,y]$ ist ein UFD, aber kein HIR für jeden Körper $k$.
Es gibt HIR, die keine euklidischen Ringe sind (z.B. $\mathbb{Z}[(1+\sqrt{-19})/2]$).
\end{proof}

\begin{example}[Euklidische Ringe]
\begin{enumerate}[label=(\alph*)]
  \item $\mathbb{Z}$ mit $\delta(n) = |n|$;
  \item $k[x]$ für jeden Körper $k$ mit $\delta(f) = \deg(f)$;
  \item Die Gaußschen ganzen Zahlen $\mathbb{Z}[i]$ mit $\delta(a+bi) = a^2 + b^2$
        (der Norm). Der Euklidische Algorithmus erlaubt die explizite Berechnung von
        $\ggT(m,n)$ in $\mathbb{Z}[i]$.
        Zwei-Quadrate-Satz: $p = a^2 + b^2$ für eine Primzahl $p$ genau dann, wenn
        $p = 2$ oder $p \equiv 1 \pmod 4$.
\end{enumerate}
\end{example}

\begin{theorem}[Eindeutige Faktorisierung in {$\mathbb{Z}[i\,]$}]
Die Gaußschen ganzen Zahlen bilden einen euklidischen Ring (also HIR und UFD).
Die Einheiten sind $\{\pm 1, \pm i\}$. Die Gaußschen Primzahlen sind:
\begin{itemize}
  \item $(1+i)$ (über $2$);
  \item $p \in \mathbb{Z}$ prim mit $p \equiv 3 \pmod 4$ (bleiben prim in $\mathbb{Z}[i]$);
  \item $a + bi$ mit $a^2 + b^2 = p$ für Primzahlen $p \equiv 1 \pmod 4$.
\end{itemize}
\end{theorem}

% ============================================================
\section{Noethersche und artinsche Ringe}
% ============================================================

\begin{definition}[Kettenbedingungen]
Ein kommutativer Ring $R$ heißt \emph{noethersch}, wenn er die \emph{aufsteigende
Kettenbedingung} (ACC, ascending chain condition) für Ideale erfüllt: Jede Kette
$I_1 \subseteq I_2 \subseteq \cdots$ von Idealen stabilisiert schließlich
($I_n = I_{n+1} = \cdots$ für ein $n$).

Äquivalent: $R$ ist noethersch genau dann, wenn jedes Ideal endlich erzeugt ist.

$R$ heißt \emph{artinsch}, wenn er die \emph{absteigende Kettenbedingung} (DCC)
erfüllt: Jede Kette $I_1 \supseteq I_2 \supseteq \cdots$ stabilisiert.
\end{definition}

\begin{theorem}[Hilbert-Basissatz]\label{satz:hilbert-basis}
Ist $R$ ein noetherscher Ring, so ist auch der Polynomring $R[x]$ noethersch.
\end{theorem}

\begin{proof}
Sei $I \subseteq R[x]$ ein Ideal. Wir zeigen, dass $I$ endlich erzeugt ist.
Für jedes $n \geq 0$ sei $J_n = \{a \in R \mid \text{es gibt } f \in I \text{ vom Grad } n
\text{ mit führendem Koeffizienten } a\} \cup \{0\}$.
Jedes $J_n$ ist ein Ideal von $R$, und $J_0 \subseteq J_1 \subseteq J_2 \subseteq \cdots$
(Multiplikation mit $x$ erhöht den Grad). Da $R$ noethersch ist, stabilisiert diese Kette:
$J_N = J_{N+1} = \cdots$ für ein $N$. Jedes $J_n$ ist endlich erzeugt: Wähle Erzeuger
$a_{n,1},\ldots,a_{n,k_n} \in J_n$ mit zugehörigen Polynomen $f_{n,j} \in I$ vom Grad $n$
mit führendem Koeffizienten $a_{n,j}$. Die endliche Menge $\{f_{n,j}\}_{n \leq N, j \leq k_n}$
erzeugt $I$ (durch ein Gradinduk{}tionsargument). Also ist $I$ endlich erzeugt.
\end{proof}

\begin{corollary}
Für jeden Körper $k$ ist der Polynomring $k[x_1,\ldots,x_n]$ noethersch.
Insbesondere hat jedes Ideal in $k[x_1,\ldots,x_n]$ ein endliches Erzeugendensystem
(eine \emph{Gröbner-Basis}, siehe Abschnitt~\ref{abschnitt:groebner}).
\end{corollary}

\begin{theorem}[Nakayama-Lemma]\label{satz:nakayama}
Sei $R$ ein lokaler Ring mit maximalem Ideal $\mathfrak{m}$, und sei $M$ ein
endlich erzeugter $R$-Modul. Wenn $M = \mathfrak{m}M$, dann $M = 0$.

Allgemeiner: Ist $N \subseteq M$ ein Untermodul und $M = N + \mathfrak{m}M$,
dann $M = N$.
\end{theorem}

\begin{proof}
Seien $m_1,\ldots,m_k$ Erzeuger von $M$. Da $M = \mathfrak{m}M$, können wir schreiben
$m_i = \sum_{j=1}^k a_{ij} m_j$ mit $a_{ij} \in \mathfrak{m}$.
In Matrixform ergibt sich $(I - A)\mathbf{m} = 0$ mit $A = (a_{ij})$.
Die Matrix $I - A$ hat Einträge $\delta_{ij} - a_{ij}$; da $a_{ij} \in \mathfrak{m}$,
sind die Diagonaleinträge $1 - a_{ii}$ Einheiten (die einzigen Nichteinheiten eines
lokalen Rings liegen in $\mathfrak{m}$). Also ist $\det(I-A)$ eine Einheit,
$I-A$ ist invertierbar, und es folgt $\mathbf{m} = 0$.
\end{proof}

\begin{corollary}
Sei $R$ lokal mit maximalem Ideal $\mathfrak{m}$ und $M$ endlich erzeugt.
Dann erzeugen $m_1,\ldots,m_k \in M$ den Modul $M$ genau dann, wenn ihre Bilder
in $M/\mathfrak{m}M$ den $R/\mathfrak{m}$-Vektorraum $M/\mathfrak{m}M$ erzeugen.
\end{corollary}

\begin{theorem}[Hopkins--Levitzki]
Ein noetherscher artinscher Ring $R$ besitzt eine Kompositionsreihe als Modul über
sich selbst. Ferner gilt: $R$ ist artinsch genau dann, wenn $R$ noethersch und von
Krull-Dimension $0$ ist.
\end{theorem}

Die \emph{Krull-Dimension} $\dim(R)$ eines Rings $R$ ist das Supremum der Längen
von Ketten von Primidealen $\mathfrak{p}_0 \subsetneq \mathfrak{p}_1 \subsetneq \cdots
\subsetneq \mathfrak{p}_n$. Für Körper gilt $\dim = 0$, für $\mathbb{Z}$ und $k[x]$
gilt $\dim = 1$, und für $k[x_1,\ldots,x_n]$ gilt $\dim = n$.

% ============================================================
\section{Lokalisierung}
% ============================================================

\begin{definition}[Lokalisierung]
Sei $R$ ein kommutativer Ring und $S \subseteq R$ eine \emph{multiplikative Menge}
(abgeschlossen unter Multiplikation, $1 \in S$). Die \emph{Lokalisierung}
$S^{-1}R$ ist die Menge der Äquivalenzklassen von Paaren $(a, s) \in R \times S$,
wobei $(a,s) \sim (a',s')$ genau dann gilt, wenn $\exists t \in S: t(as' - a's) = 0$,
mit den Operationen
\[
  \frac{a}{s} + \frac{a'}{s'} = \frac{as' + a's}{ss'}, \qquad
  \frac{a}{s} \cdot \frac{a'}{s'} = \frac{aa'}{ss'}.
\]
Es gibt einen natürlichen Ringhomomorphismus $\iota\colon R \to S^{-1}R$, $r \mapsto r/1$.
\end{definition}

\begin{example}[Wichtige Spezialfälle]
\begin{enumerate}[label=(\alph*)]
  \item $S = R \setminus \{0\}$: $S^{-1}R = \Frac(R)$ ist der Quotientenkörper eines
        Integritätsbereichs.
  \item $S = R \setminus \mathfrak{p}$ für ein Primideal $\mathfrak{p}$:
        $R_\mathfrak{p} = S^{-1}R$ ist die \emph{Lokalisierung bei} $\mathfrak{p}$,
        ein lokaler Ring mit maximalem Ideal $\mathfrak{p}R_\mathfrak{p}$.
  \item $S = \{1, f, f^2, \ldots\}$ für $f \in R$: $R_f = S^{-1}R$ invertiert $f$.
\end{enumerate}
\end{example}

\begin{theorem}[Universelle Eigenschaft der Lokalisierung]
Die Abbildung $\iota\colon R \to S^{-1}R$ ist initial unter den Ringhomomorphismen
$\varphi\colon R \to T$, für die $\varphi(s)$ eine Einheit in $T$ für alle $s \in S$
ist. Das heißt, jedes solche $\varphi$ faktorisiert eindeutig über $\iota$.
\end{theorem}

\begin{theorem}[Ideale in Lokalisierungen]
Es gibt eine Bijektion zwischen:
\begin{itemize}
  \item Primidealen von $S^{-1}R$, und
  \item Primidealen $\mathfrak{p}$ von $R$ mit $\mathfrak{p} \cap S = \emptyset$.
\end{itemize}
Unter dieser Bijektion entspricht $\mathfrak{p}$ dem Ideal $S^{-1}\mathfrak{p}$.
\end{theorem}

\begin{definition}[Dedekind-Ring]
Ein Integritätsbereich $R$ heißt \emph{Dedekind-Ring}, wenn:
\begin{enumerate}[label=(\roman*)]
  \item $R$ noethersch ist,
  \item $R$ im Quotientenkörper $\Frac(R)$ ganz abgeschlossen ist,
  \item jedes von Null verschiedene Primideal von $R$ maximal ist ($\dim R \leq 1$).
\end{enumerate}
\end{definition}

\begin{theorem}[Eindeutige Idealzerlegung in Dedekind-Ringen]
In einem Dedekind-Ring $R$ zerfällt jedes von Null verschiedene Ideal eindeutig
(bis auf Reihenfolge) in Primideale:
\[
  I = \mathfrak{p}_1^{e_1} \cdots \mathfrak{p}_r^{e_r}.
\]
\end{theorem}

\begin{proof}[Beweissk{}izze]
Die wesentlichen Schritte: (1) Jedes Ideal enthält ein Produkt von Primidealen
(noethersche Induktion); (2) für jedes Primideal $\mathfrak{p}$ ist die Lokalisierung
$R_\mathfrak{p}$ ein DVR (diskreter Bewertungsring), also HIR mit eindeutigem maximalem
Ideal; (3) die Zerlegung wird aus den lokalen Daten zusammengesetzt. Siehe \cite{Neukirch1999}.
\end{proof}

\begin{example}[Ringe ganzer Zahlen als Dedekind-Ringe]
Der Ganzheitsring $\mathcal{O}_K$ eines Zahlkörpers $K = \mathbb{Q}(\sqrt{d})$ ist
ein Dedekind-Ring. Sein ganzer Abschluss in $K$ ist:
\[
  \mathcal{O}_K = \begin{cases}
    \mathbb{Z}\!\left[\tfrac{1+\sqrt{d}}{2}\right] & \text{wenn } d \equiv 1 \pmod 4, \\
    \mathbb{Z}[\sqrt{d}] & \text{wenn } d \equiv 2 \text{ oder } 3 \pmod 4.
  \end{cases}
\]
Die \emph{Klassengruppe} $\Cl(\mathcal{O}_K)$ misst, wie weit $\mathcal{O}_K$
von einem HIR entfernt ist. Es gilt $h(d) = |\Cl(\mathcal{O}_K)| = 1$ genau dann,
wenn $\mathcal{O}_K$ ein HIR ist.
Das Stark--Heegner-Theorem klassifiziert alle imaginär-quadratischen Zahlkörper mit
$h = 1$: $d \in \{-1, -2, -3, -7, -11, -19, -43, -67, -163\}$.
\end{example}

% ============================================================
\section{Noether-Normalisierung und ganze Erweiterungen}
% ============================================================

\begin{definition}[Ganzes Element]
Sei $R \subseteq S$ eine Ringerweiterung. Ein Element $s \in S$ heißt \emph{ganz
über $R$}, wenn es eine normierte Polynomgleichung $s^n + a_{n-1}s^{n-1} + \cdots +
a_0 = 0$ mit $a_i \in R$ erfüllt. Der \emph{ganze Abschluss} von $R$ in $S$ ist der
Unterring $\bar{R} = \{s \in S \mid s \text{ ganz über } R\}$.
$R$ heißt \emph{ganz abgeschlossen}, wenn $\bar{R} = R$ (in $\Frac(R)$).
\end{definition}

\begin{theorem}[Noether-Normalisierungslemma]
Sei $k$ ein Körper und $A = k[x_1,\ldots,x_n]/I$ eine endlich erzeugte $k$-Algebra
($I$ ein Ideal). Dann gibt es algebraisch unabhängige Elemente $y_1,\ldots,y_d \in A$,
sodass $A$ eine endliche (insbesondere ganze) Erweiterung von $k[y_1,\ldots,y_d]$ ist.
Dabei ist $d = \dim A$ die Krull-Dimension von $A$.
\end{theorem}

\begin{proof}[Beweissk{}izze]
Per Induktion über $n$. Ist $I = 0$, so gilt $A = k[x_1,\ldots,x_n]$ und $y_i = x_i$.
Sonst wähle $f \in I$, $f \neq 0$. Durch eine Variablensubstitution $x_i' = x_i -
x_1^{c^{i-1}}$ (für geeignetes $c$) wird $f$ in $x_1$ nach Substitution normiert.
Dann ist $x_1$ ganz über dem kleineren Ring $k[x_2',\ldots,x_n']$, und man fährt
induktiv fort. Siehe \cite[Satz~5.10]{AtiyahMacdonald1969}.
\end{proof}

\begin{corollary}
Für eine affine Varietät $V \subseteq \mathbb{A}^n_k$, definiert durch ein Ideal
$I \subseteq k[x_1,\ldots,x_n]$, gilt:
\[
  \dim V = \dim k[V] = \dim k[x_1,\ldots,x_n]/I.
\]
Insbesondere: $\dim \mathbb{A}^n_k = n$ und $\dim V(f) = n-1$ für eine Hyperfläche.
\end{corollary}

% ============================================================
\section{Gröbner-Basen}\label{abschnitt:groebner}
% ============================================================

Gröbner-Basen, eingeführt von Bruno Buchberger im Jahr 1965~\cite{Buchberger1970},
liefern eine kanonische Menge von Erzeugern für Polynomideale und ermöglichen die
algorithmische Lösung von Ideal-Zugehörigkeits- und algebraischen Varietätsproblemen.

\begin{definition}[Monomialordnung]
Eine \emph{Monomialordnung} auf $k[x_1,\ldots,x_n]$ ist eine Totalordnung $\prec$
auf Monomen $x^\alpha = x_1^{\alpha_1}\cdots x_n^{\alpha_n}$, die mit der Multiplikation
verträglich ist und $1 \prec x^\alpha$ für alle $\alpha \neq 0$ erfüllt.

Standardbeispiele: \emph{lexikographisch} (lex), \emph{total graduiert invers
lexikographisch} (grevlex), \emph{total graduiert lexikographisch} (grlex).
\end{definition}

\begin{definition}[Divisionsalgorithmus und Gröbner-Basis]
Sei $\prec$ eine Monomialordnung fixiert. Für $f \in k[x_1,\ldots,x_n]$ bezeichne
$\operatorname{LT}(f)$ den \emph{Leitterm} (höchstes Monom mal Koeffizient).
Eine Menge $G = \{g_1,\ldots,g_s\} \subseteq I$ ist eine \emph{Gröbner-Basis} für
das Ideal $I$, wenn
\[
  \langle \operatorname{LT}(g_1), \ldots, \operatorname{LT}(g_s) \rangle
  = \langle \operatorname{LT}(f) \mid f \in I \rangle.
\]
Äquivalent: Jedes Element von $I$ reduziert sich modulo $G$ auf $0$.
\end{definition}

\begin{theorem}[Buchberger-Kriterium]
$G = \{g_1,\ldots,g_s\}$ ist genau dann eine Gröbner-Basis für $I = (G)$, wenn
für alle $i \neq j$ das \emph{S-Polynom}
\[
  S(g_i, g_j) = \frac{\operatorname{kgV}(\operatorname{LM}(g_i), \operatorname{LM}(g_j))}{\operatorname{LT}(g_i)} g_i
  - \frac{\operatorname{kgV}(\operatorname{LM}(g_i), \operatorname{LM}(g_j))}{\operatorname{LT}(g_j)} g_j
\]
sich modulo $G$ auf $0$ reduziert.
\end{theorem}

\begin{algorithm}[Buchberger-Algorithmus]
\textbf{Eingabe:} $F = \{f_1,\ldots,f_s\}$; \textbf{Ausgabe:} Gröbner-Basis $G$ für $(F)$.
\begin{enumerate}
  \item Setze $G := F$.
  \item Für alle Paare $\{g_i, g_j\} \subseteq G$ berechne $r = \overline{S(g_i,g_j)}^G$
        (Rest bei der Division durch $G$).
  \item Ist $r \neq 0$, füge $r$ zu $G$ hinzu und wiederhole.
  \item Gib $G$ aus.
\end{enumerate}
Der Algorithmus terminiert, weil das Leittermideal bei jedem Schritt streng wächst
und $k[x_1,\ldots,x_n]$ noethersch ist.
\end{algorithm}

\begin{theorem}[Anwendungen von Gröbner-Basen]
Sei $G$ eine Gröbner-Basis für $I$:
\begin{enumerate}[label=(\roman*)]
  \item \textbf{Ideal-Zugehörigkeit:} $f \in I$ genau dann, wenn $\overline{f}^G = 0$.
  \item \textbf{Lösung von Gleichungssystemen:} Elimination von Variablen durch
        Gröbner-Basen mit Lex-Ordnung; die letzte Variable erfüllt ein univariates Polynom.
  \item \textbf{Dimension:} $\dim k[x_1,\ldots,x_n]/I$ entspricht der Krull-Dimension,
        berechnet aus dem Leittermideal.
  \item \textbf{Hilbert-Funktion:} Die Anzahl der Monome, die nicht in
        $\langle \operatorname{LT}(I) \rangle$ liegen, vom Grad $\leq d$ ergibt die
        Hilbert-Funktion von $k[x_1,\ldots,x_n]/I$.
\end{enumerate}
\end{theorem}

\begin{example}
Betrachte $I = (x^2 + y - 1,\; y^2 - x) \subseteq \mathbb{Q}[x,y]$ mit Lex-Ordnung
$x \succ y$. Die Gröbner-Basis hat die Form (schematisch):
\[
  G = \{y^4 + y^2 - y - 1,\; x - y^2\}.
\]
Die Lex-Basis zeigt, dass $y$ ein univariates Polynom erfüllt, was die Lösung des
Systems $\{x^2 + y = 1, y^2 = x\}$ durch Rücksubstitution ermöglicht.
\end{example}

\begin{remark}
Die \emph{reduzierte Gröbner-Basis} ist eindeutig bestimmt (bei gegebener Monomialordnung
und $I \neq 0$): Sie ist die Gröbner-Basis, bei der alle Koeffizienten des Leitmonoms
gleich $1$ sind und kein Monom in $f_i$ durch ein Leitmonom eines anderen $f_j$
teilbar ist. Dies macht Gröbner-Basen zum kanonischen Erzeugendensystem für Polynomideale.
\end{remark}

% ============================================================
\section{Literatur}
% ============================================================

\begin{thebibliography}{10}

\bibitem{AtiyahMacdonald1969}
M.~F. Atiyah und I.~G. Macdonald.
\newblock \textit{Introduction to Commutative Algebra}.
\newblock Addison-Wesley, Reading, MA, 1969.
\newblock Das klassische prägnante Lehrbuch; behandelt Lokalisierung, noethersche
Ringe und Dimensionstheorie.

\bibitem{Lang2002}
S.~Lang.
\newblock \textit{Algebra}, überarbeitete 3.~Aufl.
\newblock Springer, New York, 2002.
\newblock Umfassendes Referenzwerk; Gaußsches Lemma, Eisenstein-Kriterium, UFDs
und Körpertheorie.

\bibitem{Matsumura1989}
H.~Matsumura.
\newblock \textit{Commutative Ring Theory}.
\newblock Cambridge Studies in Advanced Mathematics, Bd.~8.
\newblock Cambridge University Press, Cambridge, 1989.
\newblock Fortgeschrittene Behandlung noetherscher/artinscher Ringe, Dimension, Flachheit.

\bibitem{Neukirch1999}
J.~Neukirch.
\newblock \textit{Algebraische Zahlentheorie}.
\newblock Grundlehren der Mathematischen Wissenschaften, Bd.~322.
\newblock Springer, Berlin, 1999.
\newblock Dedekind-Ringe, Idealzerlegung, Klassengruppen aus zahlentheoretischer Sicht.

\bibitem{Eisenbud1995}
D.~Eisenbud.
\newblock \textit{Commutative Algebra with a View toward Algebraic Geometry}.
\newblock Graduate Texts in Mathematics, Bd.~150.
\newblock Springer, New York, 1995.
\newblock Gründliche moderne Behandlung, die kommutative Algebra und Geometrie verbindet;
Gröbner-Basen in Kapitel~15.

\bibitem{CoxLittleOShea2015}
D.~Cox, J.~Little und D.~O'Shea.
\newblock \textit{Ideals, Varieties, and Algorithms}, 4.~Aufl.
\newblock Undergraduate Texts in Mathematics.
\newblock Springer, New York, 2015.
\newblock Zugängliche Einführung in Gröbner-Basen, Buchberger-Algorithmus und
algorithmische algebraische Geometrie.

\bibitem{Buchberger1970}
B.~Buchberger.
\newblock Ein algorithmisches Kriterium für die Lösbarkeit eines algebraischen
Gleichungssystems.
\newblock \textit{Aequationes Mathematicae}, 4(3):374--383, 1970.
\newblock Originalartikel, der den Buchberger-Algorithmus einführt.

\bibitem{Noether1921}
E.~Noether.
\newblock Idealtheorie in Ringbereichen.
\newblock \textit{Mathematische Annalen}, 83(1--2):24--66, 1921.
\newblock Grundlegendes Paper, das die axiomatische Idealtheorie und die aufsteigende
Kettenbedingung etabliert.

\bibitem{Hungerford1980}
T.~W. Hungerford.
\newblock \textit{Algebra}.
\newblock Graduate Texts in Mathematics, Bd.~73.
\newblock Springer, New York, 1980.
\newblock Gute Referenz für Isomorphiesätze, Polynomringe und Moduln.

\bibitem{Dummit2004}
D.~S. Dummit und R.~M. Foote.
\newblock \textit{Abstract Algebra}, 3.~Aufl.
\newblock Wiley, Hoboken, NJ, 2004.
\newblock Standard-Graduiertenhandbuch; umfassende Behandlung der Ringtheorie
einschließlich noetherscher Ringe und UFDs.

\end{thebibliography}

\end{document}
